\documentclass[aspectratio=169,11pt]{ctexbeamer}
\usetheme{default}
\usefonttheme{professionalfonts}
\usepackage{amsmath,amssymb,amsthm,mathtools,bm}
\usepackage{booktabs,graphicx,tikz}
\definecolor{MathCatTeal}{HTML}{1F6F78}
\definecolor{MathCatMint}{HTML}{DCEFE8}
\definecolor{MathCatInk}{HTML}{18343A}
\setbeamercolor{structure}{fg=MathCatTeal}
\setbeamercolor{frametitle}{fg=MathCatInk,bg=MathCatMint}
\setbeamercolor{block title}{fg=white,bg=MathCatTeal}
\setbeamercolor{block body}{fg=MathCatInk,bg=MathCatMint!45}
\setbeamertemplate{navigation symbols}{}
\setbeamertemplate{footline}{%
  \leavevmode\hbox{%
    \begin{beamercolorbox}[wd=.82\paperwidth,ht=2.5ex,dp=1.1ex,leftskip=1em]{author in head/foot}%
      \insertshorttitle
    \end{beamercolorbox}%
    \begin{beamercolorbox}[wd=.18\paperwidth,ht=2.5ex,dp=1.1ex,center]{date in head/foot}%
      \insertframenumber/\inserttotalframenumber
    \end{beamercolorbox}%
  }%
}
\title{数学报告标题}
\author{报告人}
\institute{机构}
\date{\today}
\begin{document}
\begin{frame}
  \titlepage
\end{frame}
\begin{frame}{问题与动机}
  \begin{itemize}
    \item 用一句话陈述研究对象与问题。
    \item 说明数学动机、已知结果与报告边界。
  \end{itemize}
\end{frame}
\begin{frame}{主结果}
  \begin{theorem}
    在完整假设下陈述来源材料中的主定理。
  \end{theorem}
  \begin{alertblock}{来源纪律}
    本页的假设、量词与证明状态必须和来源材料一致。
  \end{alertblock}
\end{frame}
\begin{frame}{总结}
  \begin{enumerate}
    \item 回顾主线与关键工具。
    \item 列出限制、开放问题与引用。
  \end{enumerate}
\end{frame}
\end{document}
